%----------------------------------------------------------%
%             PARABOLA COME LUOGO E PARAMETRI              %
%----------------------------------------------------------%

\begin{question}[ID=1.1, subtitle={Parabola $y=x^2$: vertice, asse, fuoco e direttrice}, difficulty=facile, tags={parabola, luogo}]
Considera la parabola di equazione $y=x^2$. Determina vertice, asse, concavità, fuoco e direttrice e rappresentala nel piano cartesiano.
\begin{solution}$V(0,0),\ x=0,\ F\!\left(0,\tfrac{1}{4}\right),\ y=-\tfrac{1}{4}$\end{solution}
\end{question}

%----------------------------------------------------------%

\begin{question}[ID=1.2, subtitle={Parabola $x=2y^2-3y+8$: elementi caratteristici}, difficulty=facile, tags={parabola}]
Data la parabola $x=2y^2-3y+8$, determina vertice, asse, fuoco e direttrice e rappresentala nel piano cartesiano.
\begin{solution}$V\!\left(\tfrac{31}{8},\tfrac{3}{4}\right),\ y=\tfrac{3}{4},\ F\!\left(\tfrac{33}{8},\tfrac{3}{4}\right),\ x=\tfrac{29}{8}$\end{solution}
\end{question}

%----------------------------------------------------------%

\begin{question}[ID=1.3, subtitle={Parabola $(x-2)^2=12(y+1)$: elementi caratteristici}, difficulty=media, tags={parabola}]
Data la parabola $(x-2)^2=12(y+1)$, determina vertice, asse, fuoco, direttrice e concavità e rappresentala nel piano cartesiano.
\begin{solution}$V(2,-1),\ x=2,\ F(2,2),\ y=-4$, concavità $\uparrow$\end{solution}
\end{question}

%----------------------------------------------------------%

\begin{question}[ID=1.4, subtitle={Luogo equidistante da $F(0,2)$ e dalla retta $y=-2$}, difficulty=media, tags={parabola, luogo}]
Determina il luogo dei punti equidistanti da $F(0,2)$ e dalla retta $y=-2$, trovandone equazione, vertice, asse e fuoco.
\begin{solution}$x^2=8y,\ V(0,0),\ x=0,\ F(0,2)$\end{solution}
\end{question}

%----------------------------------------------------------%

\begin{question}[ID=1.5, subtitle={Parabola $y=-x^2+4x-3$: forma canonica e fuoco}, difficulty=media, tags={parabola}]
Data la parabola $y=-x^2+4x-3$, determinane vertice, asse, forma canonica, fuoco e direttrice.
\begin{solution}$y=-(x-2)^2+1,\ V(2,1),\ x=2,\ F\!\left(2,\tfrac{3}{4}\right),\ y=\tfrac{5}{4}$\end{solution}
\end{question}

%----------------------------------------------------------%

\begin{question}[ID=1.6, subtitle={Parabola con vertice $V(1,0)$ e fuoco $F(3,0)$}, difficulty=media, tags={parabola}]
Determina l'equazione e la direttrice della parabola avente vertice $V(1,0)$ e fuoco $F(3,0)$.
\begin{solution}$y^2=8(x-1),\ x=-1$\end{solution}
\end{question}

%----------------------------------------------------------%

\begin{question}[ID=1.7, subtitle={Parabola con direttrice $y=-3$ e fuoco $F(0,3)$}, difficulty=media, tags={parabola}]
Trova vertice, equazione e concavità della parabola avente direttrice $y=-3$ e fuoco $F(0,3)$.
\begin{solution}$V(0,0),\ x^2=12y$, concavità $\uparrow$\end{solution}
\end{question}

%----------------------------------------------------------%

\begin{question}[ID=1.8, subtitle={Valore di $k$ affinché la parabola $y=x^2+k$ abbia direttrice $y=1$}, difficulty=difficile, tags={parabola, parametrico}]
Determina per quali valori di $k$ la parabola $y=x^2+k$ ha direttrice $y=1$ e trova il fuoco.
\begin{solution}$k=\tfrac{5}{4},\ F\!\left(0,\tfrac{3}{2}\right)$\end{solution}
\end{question}

%----------------------------------------------------------%

\begin{question}[ID=1.9, subtitle={Luogo equidistante da $F(2,0)$ e dalla retta $x=-2$}, difficulty=difficile, tags={parabola, luogo}]
Determina il luogo dei punti equidistanti da $F(2,0)$ e dalla retta $x=-2$, trovandone equazione, vertice, asse e direttrice.
\begin{solution}$y^2=8x,\ V(0,0),\ y=0,\ x=-2$\end{solution}
\end{question}

%----------------------------------------------------------%

\begin{question}[ID=1.10, subtitle={Parabola verticale per l'origine passante per $P(2,1)$}, difficulty=difficile, tags={parabola}]
Determina equazione, fuoco e direttrice della parabola con asse verticale, vertice nell'origine e passante per $P(2,1)$.
\begin{solution}$y=\tfrac{1}{4}x^2,\ F(0,1),\ y=-1$\end{solution}
\end{question}

%----------------------------------------------------------%
%        POSIZIONE RETTA-PARABOLA                          %
%----------------------------------------------------------%

\begin{question}[ID=2.1, subtitle={Posizione di $y=x^2-4$ e $y=-2x+4$ con corda}, difficulty=facile, tags={parabola, retta, corda}]
Data la parabola $y=x^2-4$ e la retta $y=-2x+4$, determina la posizione reciproca, i punti di intersezione e la lunghezza della corda individuata.
\begin{solution}Secante; $P(2,0)$ e $P'(-4,12)$; $L=6\sqrt{5}$\end{solution}
\end{question}

%----------------------------------------------------------%

\begin{question}[ID=2.2, subtitle={Posizione di $y=x^2-2x+1$ e $y=-x+1$ con corda}, difficulty=facile, tags={parabola, retta, corda}]
Data la parabola $y=x^2-2x+1$ e la retta $y=-x+1$, determina la posizione reciproca e la corda intercettata.
\begin{solution}Secante; $L=\sqrt{2}$\end{solution}
\end{question}

%----------------------------------------------------------%

\begin{question}[ID=2.3, subtitle={Tangenza tra $y=x^2+6x+9$ e l'asse $x$}, difficulty=media, tags={parabola, tangenza}]
Studia la posizione tra $y=x^2+6x+9$ e $y=0$ determinando i punti comuni.
\begin{solution}Tangente; $P(-3,0)$\end{solution}
\end{question}

%----------------------------------------------------------%

\begin{question}[ID=2.4, subtitle={Posizione di $y=x^2$ e $y=mx+1$ al variare di $m$}, difficulty=media, tags={parabola, retta, parametrico}]
Studia la posizione tra $y=x^2$ e $y=mx+1$, determinando eventuali condizioni di tangenza.
\begin{solution}Sempre secante per ogni $m\in\mathbb{R}$\end{solution}
\end{question}

%----------------------------------------------------------%

\begin{question}[ID=2.5, subtitle={Rette orizzontali con corda di lunghezza 6 su $y=x^2$}, difficulty=media, tags={parabola, corda}]
Trova le rette orizzontali che staccano dalla parabola $y=x^2$ una corda lunga $6$.
\begin{solution}$y=9$\end{solution}
\end{question}

%----------------------------------------------------------%

\begin{question}[ID=2.6, subtitle={Valori di $k$ per intersezione di $y=2x+k$ con $y=x^2-1$}, difficulty=media, tags={parabola, parametrico}]
Determina per quali valori di $k$ la retta $y=2x+k$ ha almeno un punto in comune con la parabola $y=x^2-1$.
\begin{solution}$k\ge -2$\end{solution}
\end{question}

%----------------------------------------------------------%

\begin{question}[ID=2.7, subtitle={Valori di $m$ per cui $y=mx$ è tangente a $y=x^2+1$}, difficulty=media, tags={parabola, parametrico, tangenza}]
Determina per quali $m$ la retta $y=mx$ è tangente alla parabola $y=x^2+1$.
\begin{solution}$m=\pm 2$\end{solution}
\end{question}

%----------------------------------------------------------%

\begin{question}[ID=2.8, subtitle={Retta con corda di lunghezza 5 su $y=x^2-4x$}, difficulty=difficile, tags={parabola, corda}]
Determina l'equazione della retta parallela all'asse $x$ che stacca sulla parabola $y=x^2-4x$ una corda lunga $5$.
\begin{solution}$y=\dfrac{9}{4}$\end{solution}
\end{question}

%----------------------------------------------------------%

\begin{question}[ID=2.9, subtitle={Posizione di $y=x^2$ e $y=mx+2$ al variare di $m$}, difficulty=difficile, tags={parabola, parametrico}]
Studia la posizione tra $y=x^2$ e $y=mx+2$.
\begin{solution}Sempre secante per ogni $m\in\mathbb{R}$\end{solution}
\end{question}

%----------------------------------------------------------%

\begin{question}[ID=2.10, subtitle={Valore di $k$ affinché $y=x+k$ stacchi corda lunga 3 su $y=x^2$}, difficulty=difficile, tags={parabola, corda, parametrico}]
Trova $k$ affinché $y=x+k$ stacchi una corda lunga $3$ sulla parabola $y=x^2$.
\begin{solution}$k=\dfrac{7}{8}$\end{solution}
\end{question}

%----------------------------------------------------------%
%        RETTA TANGENTE ALLA PARABOLA                      %
%----------------------------------------------------------%

\begin{question}[ID=3.1, subtitle={Tangenti a $y=-x^2+3x$ per il punto $P(0,1)$}, difficulty=facile, tags={parabola, tangente}]
Determina le rette tangenti alla parabola $y=-x^2+3x$ passanti per il punto $P(0,1)$.
\begin{solution}$y=1$ e $y=-3x+1$\end{solution}
\end{question}

%----------------------------------------------------------%

\begin{question}[ID=3.2, subtitle={Tangenti a $y=-3x^2+2$ negli assi e area del triangolo}, difficulty=facile, tags={parabola, tangente}]
Determina le rette tangenti alla parabola di equazione $y = - 3x^2 + 2$ nei suoi punti d'intersezione con gli assi cartesiani e l'area del triangolo individuato da tali rette.
\begin{solution}Rette $y=-3x+2$; area $=2$\end{solution}
\end{question}

%----------------------------------------------------------%

\begin{question}[ID=3.3, subtitle={Tangente parallela a $y=2x$ alla parabola $y=x^2-3x+5$}, difficulty=media, tags={parabola, tangente}]
Scrivi l'equazione della retta tangente alla parabola di equazione $y = x^2 - 3x + 5$ e parallela alla retta $y = 2x$. Determina poi le coordinate del punto di contatto tra la retta e la parabola.
\begin{solution}$y=2x+3$; punto di tangenza $P(2,3)$\end{solution}
\end{question}

%----------------------------------------------------------%

\begin{question}[ID=3.4, subtitle={Tangente perpendicolare a $y=-2x$ alla parabola $y=-\frac12x^2+x$}, difficulty=media, tags={parabola, tangente}]
Determina la retta tangente alla parabola $y = -\dfrac{1}{2}x^2 + x$ perpendicolare a $y = -2x$.
\begin{solution}$y=\dfrac{1}{2}x+\dfrac{1}{2}$; punto di tangenza $P\!\left(1,\dfrac{1}{2}\right)$\end{solution}
\end{question}

%----------------------------------------------------------%

\begin{question}[ID=3.5, subtitle={Tangente perpendicolare a $y=-2x$ alla parabola (variante)}, difficulty=difficile, tags={parabola, tangente}]
Determina la retta tangente alla parabola $y = -\dfrac{1}{2}x^2 + x$ perpendicolare a $y = -2x$.
\begin{solution}$y=\dfrac{1}{2}x+\dfrac{1}{2}$\end{solution}
\end{question}

%----------------------------------------------------------%
%                FASCI DI PARABOLE - PROBLEMI              %
%----------------------------------------------------------%

\begin{question}[ID=5.1, subtitle={Fascio $y=x^2-2kx+k-1$: punti base e proprietà}, difficulty=media, tags={parabola, fasci}]
Considera il fascio di parabole di equazione $y=x^2-2kx+k-1$ e determina i punti base e la caratteristica del fascio, la parabola passante per $P(1,1)$, quella avente asse di equazione $x=1$ e le eventuali parabole tangenti alla retta $y=0$.
\begin{solution}Punti base $(0,-1)$ e $(1,0)$; asse $x=k$; $k=-1$; mai tangente\end{solution}
\end{question}

%----------------------------------------------------------%

\begin{question}[ID=5.2, subtitle={Fascio $y=x^2+kx-2k$: punti base e proprietà}, difficulty=media, tags={parabola, fasci}]
Considera il fascio $y=x^2+kx-2k$ e determina i punti base, la parabola passante per $P(2,0)$, quella con vertice sull'asse $y$ e le eventuali parabole tangenti alla retta $y=-1$.
\begin{solution}Punti base $(0,0)$ e $(2,0)$; passante per $P(2,0)$: impossibile; $k=0$; $k^2+8k-4=0$\end{solution}
\end{question}

%----------------------------------------------------------%

\begin{question}[ID=5.3, subtitle={Fascio $y=x^2+(k-1)x+k$: punti base e proprietà}, difficulty=media, tags={parabola, fasci}]
Considera il fascio $y=x^2+(k-1)x+k$ e determina i punti base, la parabola passante per $P(0,2)$, quella con asse parallelo all'asse $y$ e le parabole tangenti all'asse $x$.
\begin{solution}Punti base $(0,0)$ e $(-1,1)$; $k=2$; $k^2-6k+1=0$\end{solution}
\end{question}

%----------------------------------------------------------%

\begin{question}[ID=5.4, subtitle={Fascio $y=x^2-2kx+k-1$: variante con vertice sull'asse $y$}, difficulty=difficile, tags={parabola, fasci}]
Considera il fascio $y=x^2-2kx+k-1$ e determina i punti base, la parabola con vertice sull'asse $y$, quella passante per $P(1,1)$ e le parabole tangenti alla retta $y=1$.
\begin{solution}
Raccogliendo $k$: $k(1-2x)+(x^2-1-y)=0$. I \textbf{punti base} (indipendenti da $k$) si hanno per $1-2x=0$ e $y=x^2-1$, cioè l'unico punto $B\!\left(\tfrac12,-\tfrac34\right)$. \\
L'ascissa del vertice è $x_V=k$; \textbf{vertice sull'asse $y$} $\Rightarrow k=0$, parabola $y=x^2-1$. \\
\textbf{Passaggio per $P(1,1)$}: $1=1-2k+k-1\Rightarrow k=-1$, parabola $y=x^2+2x-2$. \\
\textbf{Tangenza a $y=1$}: da $x^2-2kx+k-2=0$, $\Delta=4(k^2-k+2)$; poiché $k^2-k+2$ ha discriminante $-7<0$ è sempre positivo, quindi $\Delta>0$ per ogni $k$: \emph{nessuna} parabola del fascio è tangente a $y=1$ (il vertice ha ordinata $-k^2+k-1\le-\tfrac34<1$, quindi la curva incontra sempre $y=1$ in due punti).
\end{solution}
\end{question}

%----------------------------------------------------------%

\begin{question}[ID=5.5, subtitle={Fascio $y=x^2+kx+k-2$: punti base e tangenza}, difficulty=difficile, tags={parabola, fasci, parametrico}]
Considera il fascio $y=x^2+kx+k-2$ e determina i punti base, la parabola passante per $P(1,0)$, quella avente vertice sulla bisettrice $y=x$ e le parabole tangenti alla retta $y=-2$.
\begin{solution}Punti base $(0,-2)$ e $(2,0)$; $k=\tfrac{1}{2}$; $k(k-4)=0$\end{solution}
\end{question}
